\section{The first split: transmitting a state with a clock}
\label{outer:section}

The schedule and end guard now permit the first split on the episode
code $E$. We will define its update before its tests, so that successful
choices cannot change which computation an episode performs. The update
is affine on $V$ and always has the original monoid action as its linear
part. Completeness of the tests is required outside an intrinsic domain
$D\subseteq E$; soundness is required for every successful pair on $E^*$.

The clock makes that soundness feasible. Equal tags let the two factors
verify a single consistent table certificate. If their tags differ, the
clock leaves only one possible ordered pair, and the affine translation
is chosen to realize its prescribed input and output. This choice is made
on the episode's actual clock graph, including episodes for which no split
will be available.

Throughout the section an episode is $[s,b)$, with $t=s+L$ and anchors
$q_0,q$ as in Section~\ref{episode:section}. The outer cut offsets form
$\mathcal H\subset B\mathbb Z$. All inspections obey the reading-side
and rejection conventions fixed there.

\subsection{The exceptional domain and a total affine update}

Define $Z(t)$ to mean that at least one of $\eta,\zeta$ is not constant
on
\[
 \{t+v:v\in B\mathbb Z,\ |v|\le w\}.
\]
For a position $Q$, let $u=(Q-t)\bmod B$, with residues represented in
$\{0,\ldots,B-1\}$.  Define $A_*(t,Q)$ to be the disjunction of the
following three conditions:
\begin{enumerate}
\item $\varphi(Q-t)\in W$;
\item for some $0\le j\le g$ and some integer $v$ with $|v|\le w_0$,
      $f_j^u(t+v)\ne f_j^u(t)$;
\item for some $0\le j\le g$, $f_j^u(t)=\bot$, but the inclusive
      interval
      $[z_j(t),z_j(t)+H_u+H_{\rm ext}]$ contains an inner marker.
\end{enumerate}
In the second condition the index $u$ is held fixed as $t$ is shifted.
The search values are absolute positions, with $\bot$ equal only to
$\bot$.  Set
\begin{equation}
 D=\{e\in E: Z(t)\ \text{or}\ A_*(t,q)\}.
 \label{outer:exceptional-domain}
\end{equation}
The episode determines both predicates without choosing a split.
Thus $D$ is a fixed language. In particular, failure of a particular
choice of tag is not its definition; the same $D$ must serve every input
state and every pair of independently chosen tags.

For $A,C\in M$, write
\begin{equation}
 f_{A,C}(a_1,\ldots,a_g)\colon
 v\longmapsto vA a_1\cdots a_g C.
 \label{outer:product-table}
\end{equation}
This is a map from $M^g$ into $\operatorname{End}_{\rm lin}(V)$.
For every such map, use the shift table $\beta_{f_{A,C}}$ fixed in
Lemma~\ref{alg:shift-table}.  We define $F_e\colon V\to V$ on
\emph{every} complete episode, including every episode in $D$:
\begin{enumerate}
\item If $\mathcal E_{\varphi(q-t)}$ is the singleton off-diagonal edge
      $\{(\alpha,\lambda)\}$, put
      \begin{equation}
      F_e(v)=vT_{s,b}+Y(\lambda)-X(\alpha)T_{s,b}.
      \label{outer:off-diagonal-update}
      \end{equation}
\item Otherwise put $u=(q-t)\bmod B$.  If at least one
      $f_j^u(t)$ is absent, put $F_e(v)=vT_{s,b}$.
\item In the remaining case let $l_j=f_j^u(t)$ for $0\le j\le g$.
      Their order is $l_0<\cdots<l_g$ by the schedule.  Set
      \[
      A=T_{s,l_0},\qquad C=T_{l_g,b},\qquad
      d_i=T_{l_{i-1},l_i}\quad(1\le i\le g).
      \]
      Then put
      \begin{equation}
      F_e(v)=v f_{A,C}(\mathbf d)+\beta_{f_{A,C}}(\mathbf d).
      \label{outer:table-update}
      \end{equation}
\end{enumerate}
These three cases define exactly one map on every episode. The first
case is unambiguous by the clock alternative. The other cases cover all
graphs contained in the diagonal, including the empty graph, with the
actual residue and actual marker data. In the third case,
$A d_1\cdots d_g C=T_{s,b}$, so all three formulas have the same
prescribed linear part. Hence $F_e$ is total and deterministic even on
$D$, and no cancellation or inverse in $M$ has entered its definition.

\subsection{The two independently chosen tags}

For each $x,y\in V$ we define languages $P_x,Q_y$ by finite unions
over tags.  A prefix factor is read on $[s,k)$, and a suffix factor
on $[k,b)$.  These positions describe the tests in their own relative
coordinates; they do not give either factor access to the other one.

Every prefix test chooses a tag $\alpha$, requires
$x=X(\alpha)$, chooses a cut of the form specified by that tag, and
requires
\begin{equation}
                 \varphi(k-t)\in I_\alpha.
 \label{outer:prefix-clock}
\end{equation}
Every suffix test chooses an independent tag $\lambda$.  It performs
the shared-end guard of Lemma~\ref{episode:shared-end}, using all
origins $k-a$, $a\in\mathcal H$.  More explicitly, both end searches
$\eta,\zeta$ must be constant over these origins; the suffix must
end exactly as prescribed at one fixed reference origin; all the
required seed and inspection data must be present.  Let $q$ denote
the clock anchor so obtained.  The suffix requires
\begin{equation}
 y=Y(\lambda),\qquad
 \varphi(q-k)\in J_\lambda,\qquad
 u(\lambda)=(q-k)\bmod B.
 \label{outer:suffix-clock}
\end{equation}
The field $u(\lambda)$ is part of the suffix's tag.  The additional
tests depend on the tag's form as follows.

\paragraph{Main tags.}
A main tag has fields $(j,u,A,C,\mathbf d_{-j},X,Y)$.
Its prefix chooses one of the cuts
\[
                 k=t+c_j+aB,\qquad 0\le a<R.
\]
Using the tag's fixed index $u$, it requires the boundary searches
$f_0^u(t),\ldots,f_{j-1}^u(t)$ to be present.  It checks $A$ and
the listed products $d_i$ for $i<j$ against the intervals on its
side of the cut.  For $a\in M$, let $\mathbf d[j\leftarrow a]$
denote the tuple formed by inserting $a$ into position $j$ of the
tag's listed entries.  The prefix also checks the finite table certificate
\begin{equation}
 X f_{A,C}(\mathbf d[j\leftarrow a])
 +\beta_{f_{A,C}}(\mathbf d[j\leftarrow a])=Y
 \qquad(a\in M).
 \label{outer:main-certificate}
\end{equation}
The fields used in this certificate are the tag's fields, not an
assumption about the unseen half of the episode.

A suffix with main tag $(j,u,A,C,\mathbf d_{-j},X,Y)$ requires all
searches with labels $j,\ldots,g$, with this fixed $u$, to be present
and to have identical absolute positions at every origin
\begin{equation}
                  k-c_j-aB,\qquad 0\le a<R.
 \label{outer:main-hypotheses}
\end{equation}
Using those common positions, it checks $d_i$ for $i>j$ and $C$
against the suffix.  Every inspection is required at every origin
in \eqref{outer:main-hypotheses}; an unavailable inspection rejects
the suffix instead of discarding that origin.  The reading-side
separation in the main-cut schedule ensures that these are suffix
tests even when the tag is not the tag chosen by a paired prefix.

\paragraph{Periodic tags.}
A periodic tag has fields $(j,u,m_p,X,Y)$, where $0\le j\le g$.
Its prefix chooses a cut
\[
 k=z_j(t)+H_u+L_{\rm lag}+aB,\qquad 0\le a<n d N_M,
\]
checks $f_j^u(t)=\bot$, and checks $m_p=T_{s,k}$.
A suffix with this tag requires
\begin{equation}
                       Y=X m_p T_{k,b},
 \label{outer:periodic-certificate}
\end{equation}
where $b$ is its own exact endpoint, as required by the shared-end
guard.  It does not presume that the independently chosen prefix
has the same tag.

These rules completely specify $P_x$ and $Q_y$.  All choices of
tags, cut offsets, table entries, and origins belong to previously
fixed finite sets.  In particular, the universal certificate in
\eqref{outer:main-certificate} is a finite Boolean test.

\subsection{Correctness for arbitrary successful pairs}

\begin{proposition}[First split]\label{outer:split}
For the total updates $F_e$ above and $D\subseteq E$ from
\eqref{outer:exceptional-domain}, the languages $P_x,Q_y$ have the
following properties.
\begin{enumerate}
\item On any word in $E^*$, a match of $P_xQ_y$ beginning at the
      beginning of the word consumes exactly its first episode $e$
      and satisfies $F_e(x)=y$.
\item For every $e\in E\setminus D$ and every $x\in V$, some cut
      of $e$ gives a match of $P_xQ_{F_e(x)}$.
\end{enumerate}
There is no match on the empty word.
\end{proposition}

\begin{proof}
\emph{Soundness.} Let independent tags $\alpha,\lambda$ witness a
successful pair at the beginning of a word in $E^*$. The prescribed
prefix length is positive and, on a nonempty such word, ends at an outer
cut strictly within its first episode. The true origin is among the
origins checked by the suffix. Lemma~\ref{episode:shared-end} therefore
forces the true endpoint $b$ and clock anchor $q$, even when the two tags
propose different forms or boundary indices. On the empty word the
positive prefix length cannot match.

The clock tests now add using actual common positions:
\[
 \varphi(q-t)=\varphi(k-t)+\varphi(q-k)\in I_\alpha+J_\lambda.
\]
If the graph is contained in the diagonal, this edge forces
$\alpha=\lambda$. Since $k-t$ is divisible by $B$, the suffix's residue
field is the true $u=(q-t)\bmod B$. For a main tag, the actual origin
occurs among its role-specific hypotheses
\eqref{outer:main-hypotheses}. The prefix verifies boundaries through
$j-1$, the suffix verifies those from $j$ onward, and together they
verify $A,C$ and every listed interval product. All actual boundaries
are therefore present. Substituting the actual missing product into the
universal certificate \eqref{outer:main-certificate} gives
$F_e(x)=y$ under \eqref{outer:table-update}.

For an equal periodic tag, the prefix establishes absence of the actual
search with the true index $u$. The linear case of $F_e$ therefore
applies. Its verified $m_p=T_{s,k}$ and the suffix certificate imply
\[
 y=xm_pT_{k,b}=xT_{s,b}=F_e(x).
\]
An empty clock graph cannot produce either successful case.

If the graph instead consists of one off-diagonal edge, the successful
pair must be exactly that edge. The separate state checks say
$x=X(\alpha)$ and $y=Y(\lambda)$. Formula
\eqref{outer:off-diagonal-update} was defined to give this equality
$F_e(x)=y$. No agreement of their other fields is needed. This exhausts
all successful pairs and proves soundness.

\emph{Completeness.} Fix $e\in E\setminus D$ and an arbitrary input
$x\in V$. The negation of $Z(t)$ supplies the common end guard at every
scheduled outer cut. Put $u=(q-t)\bmod B$. Since
$\varphi(q-t)\notin W$, every tag $\alpha$ has a ball of radius $\rho$
in $I_\alpha\cap(\varphi(q-t)-J_\alpha)$, and the clock graph is the
full diagonal. We must find a cut that hits the ball for a tag whose
letter tests and certificate are correct.

If all searches $f_i^u(t)$ are present, take their true products
$A,C,\mathbf d$. The shift-table lemma supplies a coordinate $j$ that
can be omitted for this tuple and input. Give the main tag those true
fields and $X=x$, $Y=F_e(x)$. Its universal certificate then holds. The
clock values at its main cuts are
\[
 \varphi(c_j)+\varphi(aB),\qquad 0\le a<R,
\]
a translate of the fixed net, so one cut meets its ball. At that cut,
every alternative role-specific origin is $t+(a-a')B$ with displacement
of magnitude less than $w_0$. The stability part of $\neg A_*(t,q)$,
with $u$ fixed, makes the required future marker positions agree.
The remaining tests are actual products on the proper reading sides.
Thus the main pair succeeds.

If some $f_j^u(t)$ is absent, write $z=z_j(t)$. The last part of
$\neg A_*(t,q)$ makes $[z,z+H_u+H_{\rm ext}]$ markerless.
Lemma~\ref{schedule:periodic-cuts} gives a period $d'<d$ and cuts
\begin{equation}
 k_a=k_0+a d'N_MB,\qquad 0\le a<n,\qquad
 k_0=z+H_u+L_{\rm lag},
 \label{outer:periodic-progression}
\end{equation}
whose prefix products share one value $m_p$. Choose the periodic tag
with these $j,u,m_p$ and with $X=x$, $Y=xT_{s,b}=F_e(x)$.
It is a single tag for the whole progression: constancy of $m_p$ follows
from eventual periodicity of monoid powers, without cancelling an earlier
factor. The same lemma shows that its clock values traverse a translate
of the full clock image, and hence include a translate of the fixed net.
One therefore hits the ball for this tag. The marker-absence inspections
are before the cut by the lag condition, and the suffix certificate
$Y=Xm_pT_{k_a,b}$ follows from $m_pT_{k_a,b}=T_{s,b}$. The common end guard already holds. This proves
completeness in the periodic case as well, for every input $x$.
\end{proof}

We now have the first split interface on all of $E$, with the same
updates $F_e$ reserved for its skipped episodes. To complete the sequence
computation, it remains to handle $D^*$. The defining exceptional
predicate, which was an obstruction to the clock construction, will
serve as a witness to the origin at an earlier stencil cut. The next
section proves that this witness is usually unique across the entire
stencil.
